\subsection{The CM punctured elliptic curve and its nilpotent quotients}
\paragraph{Statement E1 (Proved; Part A(f) and the class-two part of E(c)).}\label{res:E1}
Put $K=\mathbb Q(i)$, $E:y^2=x^3-x$, $U=E\setminus\{O\}$ and $\Delta=\Gamma_{K^{\rm ab}}$. Use a $K$-rational tangent at $O$ to define $s$. Let $P=\pi_1(U_{\overline K})^{(5)}$, let $\gamma_rP$ denote its closed lower central series, and put $N_r=P/\gamma_{r+1}P$. Then
\[
 N_2^{s(\Delta)}=N_2,\qquad
 N_2^{s(\Gamma_K^\circ)}=\gamma_2P/\gamma_3P\simeq\mathbb Z_5(1).
 \tag{E.2}
\]
The full profinite group has $C(s)=I_O$. Its image in $\pi^{\rm ab}=T(E)$ is zero, and its image in $N_2$ is exactly the centre $\gamma_2P/\gamma_3P$. Thus the first group in (E.2) contains non-inertia elements in this quotient. The comparison is with LM3b D.5 and Longue Marche section 34, scan 483.

The image assertion has the generality asked for in Part A(f): for every
elliptic curve $E_0$ over a number field $K_0$, every rational tangential
section at its removed origin, and every prime $\ell$, the image of its
full $C(s)$ in $T(E_0)$ is zero and its image in the class-two geometric
pro-$\ell$ quotient is the entire central $\mathbb Z_\ell(1)$.
In particular this applies to every CM elliptic curve over a number field
containing its CM field. The additional fixed-group formulas (E.2) retain
the specific curve, subgroup and prime stated above.

\begin{proof}
The topological punctured torus has free group on two generators $x,y$, with its boundary represented by $c=[x,y]$ up to inversion and conjugation. SGA~1 XII.5.1 identifies its pro-$5$ completion with $P$, so $P$ is free pro-$5$ of rank two. The quotient $N_2$ has unique coordinates $x^a y^b c^d$ with $a,b,d\in\mathbb Z_5$. Existence follows by collecting commutators, which are central; uniqueness follows by mapping to the upper triangular unipotent $3$ by $3$ matrices over $\mathbb Z_5$. Hence
\[
 1\longrightarrow Z=\mathbb Z_5c\longrightarrow N_2
 \longrightarrow V=T_5(E)\longrightarrow1,
 \qquad Z=\bigwedge^2V=\mathbb Z_5(1).
\]
The determinant action on $Z$ follows by computing the commutator of two lifted vectors; the last identification is the Weil pairing.

The endomorphism $[i]:(x,y)\mapsto(-x,iy)$ is defined over $K$ and commutes with Galois. Its action $J$ on $V$ has $J^2=-1$. The two roots of $X^2+1$ in $\mathbb Z_5$ differ by a unit. Their idempotents decompose $V$ into two free rank-one eigenspaces. Both have rank one because on the complex homology lattice multiplication by $i$ has characteristic polynomial $X^2+1$, and comparison preserves that integral action. Each Galois action on one rank-one summand is a character, hence factors through $\Gamma_K^{\rm ab}$. Therefore $\Delta$ acts trivially on $V$ and on $Z$.

We must also compute the possible shear from $V$ to $Z$ in the actual tangential action. The parameter $t=-x/y$ at $O$ satisfies $[-1]^*t=-t$ exactly. On the local pro-$5$ covers $t_n^{5^n}=t$ the maps $t_n\mapsto-t_n$ are compatible, square to the identity and commute with every coefficient automorphism. They define the Galois-invariant path between the two opposite tangential base points. Composing this path with $[-1]$ gives a genuine involution $\alpha$ of $P$ commuting with $s(\Gamma_K)$. This construction uses the odd prime $5$ and asserts no corresponding order-two path for the full profinite group.

Since $2$ is invertible in $\mathbb Z_5$, the class-two group law can be written in integral Campbell--Hausdorff coordinates on a Lie module $L$, with product $u*v=u+v+\tfrac12[u,v]$. The involution is $-1$ on $L/Z=V$ and $+1$ on $Z$. The idempotents $(1\pm\alpha)/2$ give $L=L^-\oplus Z$ and the projection $L^-\to V$ is an isomorphism. Every Galois element preserves this splitting. Its action is therefore exactly its action on $V$ together with the determinant action on $Z$, with no additional shear. This proves $N_2^{s(\Delta)}=N_2$.

For $\Gamma_K^\circ$, A1 gives $V^{\Gamma_K^\circ}=0$, while $Z$ is fixed because its character is $\chi_5$. The same splitting proves the second equality in (E.2). In the full group A6 gives $C(s)=I_O$. The boundary $[x,y]$ is zero in abelianization and generates $Z$ in $N_2$, proving both asserted images. A different rational tangent changes the action by a cocycle with values in $I_O$. Its image in $N_2$ is central, so all the class-two conclusions are independent of this change.

For the general image assertion, A6 again identifies the full centralizer
with origin inertia. The punctured-torus presentation and the collection
map into the Heisenberg group apply over $\mathbb Z_\ell$ for every
$\ell$. They identify its boundary commutator with a generator of the
central factor, and its abelian image with zero. These are the two actual
quotient maps used in the assertion; the general argument needs no
Campbell--Hausdorff coordinate change.
\end{proof}

\paragraph{Statement E2 (Proved; audit of LM3b D.5).}\label{res:E2}
In the class-two arithmetic pro-$5$ quotient for this curve, the cyclic subgroup generated by a lift $x$ of a primitive vector in a CM eigenspace of $V$ is maximal procyclic, has arithmetic normalizer mapping onto $\Gamma_K$, and is not a cusp inertia subgroup in that quotient.

\begin{proof}
Choose a second lift $y$ completing that vector to a basis of $V$ and write $c=[x,y]$. Every lift of $\sigma\in\Gamma_K$ acts on $x$ by
\[
 x\longmapsto x^{a_\sigma}c^{d_\sigma},\qquad
 a_\sigma\in\mathbb Z_5^*,\ d_\sigma\in\mathbb Z_5.
\]
Replacing that arithmetic lift by its product with $y^{d_\sigma/a_\sigma}$ removes the central factor, because
\[
 y^r x^a c^d y^{-r}=x^a c^{d-ar}.
\]
Thus the normalizer maps onto $\Gamma_K$. The primitive nonzero image of $x$ in $V$ proves maximality: a proper overgroup that is procyclic would express $x$ as a proper $5$-power in that group, forcing its image in $V$ to be divisible by $5$. All cusp inertia here lies in $Z$, while this subgroup has nonzero image in $V$. These conclusions use the explicit quotient $\pi\twoheadrightarrow N_2$; they do not assert a lift of its normalizer to the full arithmetic fundamental group.

The program checks every unit $a$ and every $d$ modulo $5,25,125$, respectively $20$, $500$, $12500$ conjugations. It also computes the two distinct roots of $X^2+1$ at those moduli, verifying the eigenspace denominators used in E1.
\end{proof}

\paragraph{Statement E3 (Open in the full group; finite classes proved).}\label{res:E3}
For the actual $s(\Delta)$-action above, the entire last central layer $\gamma_rP/\gamma_{r+1}P$ is fixed in $N_r$ for every $r\ge1$. For every $r\ge3$, it contains a nonidentity fixed element outside every conjugate of the image of cusp inertia. The calculations through $r=10$ give the degree ranks
\[
 2,1,2,3,6,9,18,30,56,99.
\]
They do not determine whether the full group $P^{s(\Delta)}$ contains a non-inertia element. The same unresolved issue remains for $\pi^{s(\Delta)}$ and for the character-free maximal-procyclic assertion of LM3b D.4.

\begin{proof}[Proof and exact remaining obligation]
An automorphism inducing the identity on $P/\gamma_2P$ induces the identity on every $\gamma_rP/\gamma_{r+1}P$: write generators of the latter as length-$r$ iterated commutators, replace each generator by itself times an element of $\gamma_2P$, and collect. Each error has length at least $r+1$. These commutators topologically generate the layer by the definition of the closed lower central series and the commutator identities. E1 shows that each element of $s(\Delta)$ has this property. Since the last layer is central in $N_r$, its elements are actual fixed group elements, not merely fixed vectors in a different representation.

Take $w_r=[x,[x,\ldots,[x,y]\ldots]]$ of length $r$. Under the Magnus substitution $x\mapsto1+X$, $y\mapsto1+Y$, its lowest term is $\operatorname{ad}_X^{r-1}(Y)$, whose coefficient of $X^{r-1}Y$ is $1$. Mapping into the finite group of units with constant term $1$ in
\[
 (\mathbb Z/5)[X,Y]_{\rm noncomm}/(\text{all words of degree }>r)
\]
makes every commutator of length $r+1$ trivial and keeps this coefficient nonzero. The universal property of the free pro-$5$ group gives that continuous finite quotient. This proves $w_r\ne1$ in $N_r$. Every power $c^a$ has degree-two image $a c$ in the torsion-free centre of $N_2$. Consequently the image of any conjugate of the cusp inertia intersects the last layer of $N_r$, $r\ge3$, trivially. This proves the asserted non-inertia witness for every finite class.

The program enumerates the Lyndon words in $X,Y$ through degree ten and expands their standard bracketings in the associative word algebra over $\mathbb Z$. The leading word of each expansion has coefficient one and all other terms are larger. This gives an exact triangular independence certificate. Collection using antisymmetry and the Jacobi identity expresses all Lie words in these bracketings; it is obtained inductively by factoring a Lyndon word at its longest Lyndon suffix. The counted ranks agree with $r^{-1}\sum_{d\mid r}\mu(d)2^{r/d}$. The explicit witness expansions and all ranks are in the output. The fixed-layer conclusion follows from the preceding commutator proof, rather than from treating arbitrary matrices as Galois images.

For the passage to the actual pro-$5$ graded layer, make the same
truncated-unit substitution over $\mathbb Z/5^m$ for every $m$.
The triangular leading coefficients remain one at each such level.
Consequently any $\mathbb Z_5$-linear relation between the displayed
brackets has every coefficient divisible by $5^m$ for every $m$, and
all coefficients vanish. Together with the collection argument, this
proves the stated $\mathbb Z_5$-ranks in the geometric group itself.

For the precise missing passage, the equivariant quotient $q_r:P\to N_r$ gives
\[
 q_r(P^{s(\Delta)})\subseteq N_r^{s(\Delta)},\qquad
 P^{s(\Delta)}=\varprojlim_r N_r^{s(\Delta)}.
 \tag{E.3}
\]
The second equality uses $P=\varprojlim_rN_r$: every finite $5$-group is nilpotent, so every finite quotient of $P$ kills some $\gamma_rP$, and compactness proves both injectivity and surjectivity of this inverse-limit identification. The explicit $w_r$ maps to $1$ in $N_{r-1}$. Thus these particular fixed elements do not form a nontrivial compatible element of (E.3). Already at the transition $N_3\to N_2$, lifting a specified noncentral fixed $v\in N_2$ requires vanishing of the continuous cocycle
\[
 \delta\longmapsto s(\delta)(\widetilde v)\widetilde v^{-1}
 \ \in\gamma_3P/\gamma_4P.
\]
It is independent of the chosen lift, because this kernel is fixed pointwise. Its vanishing for all $\delta\in\Delta$ is exactly the missing lift condition at this first transition; success at it still leaves the compatible conditions at every higher transition. This paper computes neither that arithmetic cocycle nor a compatible non-inertia tower. In the full profinite group there is the additional requirement of compatible lifts through its finite quotients beyond the pro-$5$ quotient. These are the explicit remaining obligations; the quotient counterexample E2 does not dispose of them.
\end{proof}
